\begin{figure}[htbp]
\centering
\begin{tikzpicture}
\begin{axis}[ybar, width=11cm, height=6cm, bar width=9pt,
  ymode=log, ymin=0.7, ymax=1500, log origin=infty,
  xlabel={time bucket (\textmu s)}, ylabel={samples (log scale)},
  symbolic x coords={{<51},{51--177},{177--304},{304--430},{557--683},{$\ge$1063}},
  xtick=data, x tick label style={font=\tiny, rotate=20, anchor=east},
  tick label style={font=\scriptsize}, label style={font=\small},
  point meta=explicit symbolic, nodes near coords, nodes near coords style={font=\tiny},
  grid=major, grid style={gray!20}]
\addplot[fill=accent!60, draw=accent] coordinates
  {({<51},500) [500] ({51--177},488) [488] ({177--304},7) [7] ({304--430},1) [1] ({557--683},3) [3] ({$\ge$1063},1) [1]};
\end{axis}
\end{tikzpicture}
\caption{The 1000 timings of one 1\,MiB sum from
\texttt{ch15\_cdf} (Linux container) as a histogram with a
logarithmic count axis; empty buckets (430--557, 683--1063) are
omitted. 98.8\% of the samples sit at or below 177\,\textmu s, but
twelve samples lie out to a millisecond. The median is 50.7\,\textmu s;
the maximum is 21 times larger.}
\label{fig:hist15}
\end{figure}
